\section{Access model and polynomial implementation}\label{sec:model}
An exact block encoding of a Hermitian contraction $H$ is a unitary $U$ and a known isometry $J\colon v\mapsto |0^a\rangle\otimes v$, with $J^*UJ=H$. A unit-normalized $\eta$-approximate block encoding of $A$ is a unitary $V$ and a known isometry $J_{\mathrm{out}}$ such that
\[
 \norm{J_{\mathrm{out}}^*VJ_{\mathrm{out}}-A}\leq\eta.
\]
An exact block encoding of $A$ with normalization $\nu>0$ instead satisfies $J_{\mathrm{out}}^*VJ_{\mathrm{out}}=A/\nu$; unit normalization means $\nu=1$.
All operator norms are spectral norms. A converter is a circuit whose gates may depend on the stated spectral promise and accuracy, but not on the unknown matrix or on the completion of its oracle. It must work for every unitary $U$ with the prescribed compressed block. One controlled call to $U$ or $U^*$ costs one query. Ancillas, inverse calls, and coherently controlled choices of queries are allowed. The output is a reusable unitary; it is not conditioned on a favorable measurement outcome. If the worst-case query count is bounded, intermediate measurements can be deferred and their classical controls made coherent, so the lower bounds below cover this form of adaptivity. We do not consider algorithms with unbounded stopping times and bounded expected cost.

The main spectral promise is
\begin{equation}\label{eq:bands}
 0<\delta,\qquad \rho>1,\qquad 0<c\leq1,\qquad
 \rho\delta<c,\qquad
 \spec(H)\subseteq[\delta,\rho\delta]\cup[c,1].
\end{equation}
Let $Q(\delta,\eta;\rho,c)$ be the minimum worst-case query count for unit-normalized conversion to $I-H$ under this promise, uniformly in the matrix dimension. Constants in asymptotic notation may depend on the fixed parameters shown as subscripts. In the fixed-relative-accuracy results, $\rho,c,K$ are fixed and $\eta=K\delta$ as $\delta\downarrow0$.

\begin{lemma}[Bounded polynomial implementation]\label{lem:implementation}
If $p\in\R[x]$ has degree $d$ and $|p(x)|\leq1$ on $[-1,1]$, then $p(H)$ has an exact unit-normalized block encoding using $O(d)$ controlled queries to any exact block encoding of the Hermitian contraction $H$. If $p$ has definite parity, $d$ queries suffice in the real-polynomial QSVT convention.
\end{lemma}
\begin{proof}
The definite-parity statement is the real-polynomial singular value transformation theorem of Gily\'en et al.\ \cite[Corollary 18, extended version]{gilyen2019}. It extracts the real part by coherently averaging the phase sequences with opposite signs; this uses an extra qubit and loses no normalization. For any Hermitian $H$, the even transform gives $p(|H|)=p(H)$, while the odd transform includes the left/right singular-vector signs and gives $\operatorname{sign}(H)p(|H|)=p(H)$. On the zero eigenspace the odd case follows from $p(0)=0$.

For completeness, we give the reduction for polynomials without definite parity; it makes no assumption on the unknown completion. Define the Hermitian unitary
\[
 \widetilde U=|0\rangle\langle1|\otimes U+
 |1\rangle\langle0|\otimes U^*,\qquad
 \widetilde J=|+\rangle\otimes J.
\]
Its compression is $H$, and it costs a constant number of controlled oracle calls. Set $\Pi=\widetilde J\widetilde J^*$ and $V=(2\Pi-I)\widetilde U$. For each eigenvector $v$ of $H$ with eigenvalue $x\in(-1,1)$, the span of $\widetilde Jv$ and $(\widetilde U\widetilde Jv-x\widetilde Jv)/\sqrt{1-x^2}$ is invariant, and $V$ has matrix
\[
 \begin{pmatrix}x&\sqrt{1-x^2}\\-\sqrt{1-x^2}&x\end{pmatrix}.
\]
Thus $(V+V^*)/2=xI$ on this subspace; the endpoint eigenspaces follow directly. The ordinary polynomial
\[
 R(z)=z^d p\bigl((z+z^{-1})/2\bigr)
\]
has degree at most $2d$ and modulus at most one on the unit circle. Generalized quantum signal processing (GQSP) \cite{motlagh2024}, via a Fej\'er--Riesz completion, implements $R(V)$ with at most $2d$ controlled walk queries. Multiplication by $V^{-d}$ uses another $d$ queries and gives $p((V+V^*)/2)$. Compression by $\widetilde J$ gives $p(H)$. The construction is coherent and unit-normalized throughout.
\end{proof}

\begin{lemma}[Scalar representation and derivative bounds]\label{lem:scalar}
For the scalar oracle
\begin{equation}\label{eq:scalar-oracle}
 U(t)=\begin{pmatrix}\sin t&\cos t\\\cos t&-\sin t\end{pmatrix},
\end{equation}
each matrix element $q(t)$ of a $T$-query converter is a trigonometric polynomial of degree at most $T$ and satisfies $|q(t)|\leq1$ for every real $t$. The function $g(x)=1-\operatorname{Re}q(\arcsin x)$ obeys $0\leq g\leq2$ on $[-1,1]$ and, for each fixed integer $j\geq1$,
\begin{equation}\label{eq:derivative-bound}
 \sup_{|x|\leq1/2}|g^{(j)}(x)|\leq C_j(T+1)^j.
\end{equation}
\end{lemma}
\begin{proof}
The entries of each query, including a controlled query whose inactive block is the identity, have trigonometric degree at most one. Each query therefore increases the degree of the entries by at most one. Gates that do not depend on the oracle, and ancillas, do not affect this induction. Unitarity bounds each output entry for all $t$, whether or not $\sin t$ satisfies the spectral promise. Trigonometric Bernstein inequalities bound the $j$th derivative in $t$ by $2T^j$ for $1-\operatorname{Re}q(t)$. Repeated differentiation of $t=\arcsin x$, whose derivatives of every fixed order are bounded on $[-1/2,1/2]$, proves \eqref{eq:derivative-bound}.
\end{proof}
